\section{Lecture 6 (January 24, 2025): Symmetric Braidings and Rigidity}\label{lec6}
\noindent\textit{Scribe: Jaedon Rich}
 

\begin{definition}
A braiding $c_{X,Y}: X \otimes Y \to Y \otimes X$ is \textit{symmetric} if $$c_{Y,X} c_{X,Y} = \Id_{X \otimes Y}.$$
\end{definition}

\begin{remark}
\hfill
    \begin{itemize}
        \item Let $\cC$ be braided and take $X \in \cC$. Then we get an action of the braid group 
        \[
        B_n=\langle b_1,\ldots,b_{n-1}\mid
        \begin{array}{ll}
        b_ib_j=b_jb_i&(|i-j|>1),\\
        b_ib_{i+1}b_i=b_{i+1}b_ib_{i+1}&(1\leq i\leq n-2)
        \end{array}\rangle
        \]
        on $X^{\otimes n}$, with $b_i$ acting by $c_{X,X}$ in positions $i,i+1$ and identities elsewhere. Coherence supplies the required associators.
        \item If the braiding is symmetric, this action factors through $S_n = B_n/\langle b_i^2 \rangle$
    \end{itemize}
\end{remark}
\begin{example}
The ``swap'' braiding in $\bVec_k$
\[x \otimes y \mapsto y \otimes x\]
is symmetric.
\end{example}
\begin{example}
For every affine group scheme $G$, the ordinary flip is $G$-equivariant for the diagonal action, so its representation category is symmetric as well.
\end{example}

\begin{example}
    Take an abstract group $G$. Define the category $\bVec_k G$ of $G$-graded $k$-vector spaces whose objects are vector spaces of the form $$V = \bigoplus_{g \in G} V_g$$ and whose morphisms are degree-preserving linear maps.\\
    For any $G$, this has a monoidal structure given by
    \[(V \otimes W)_g = \bigoplus_{xy = g} (V_x \otimes W_y).\]
    In this case, a braided structure exists if and only if $G$ is abelian. To see the only if direction, notice that for any $g \in G$ we have an object $\delta_g$ which is $k$ in degree $g$ and $0$ in all other degrees. If $G$ has a braiding, then we get
    \[\delta_{gh} = \delta_g \otimes \delta_h \simeq \delta_h \otimes \delta_g = \delta_{hg},\]
    and so $gh = hg$. Conversely, for abelian $G$, the ordinary flip preserves the grading and satisfies the hexagon identities for the usual associator.
\end{example}
\begin{example}
    Consider $\SVec_k = \bVec_k(\Z/2\Z)$, vector spaces graded by $\Z/2\Z$. In particular, every object is of the form $V = V_0 \oplus V_1$, where we call $V_0$ and $V_1$ the even and odd homogeneous components, respectively. \\
    For $v \in V_i$, we define $|v| = \deg v = i$. Then the braiding $V \otimes W \to W \otimes V$ is given on homogeneous elements by
    \[v \otimes w \mapsto (-1)^{|v||w|} w \otimes v.\]
    The sign is multiplicative in each degree, so the hexagon identities hold; applying the flip twice gives sign $(-1)^{2|v||w|}=1$, hence symmetry. When $\operatorname{char}k=2$, the sign equals $1$ and this braiding is the ordinary flip on graded spaces. The distinction considered below requires characteristic different from $2$.
\end{example}
\begin{definition}
    Fix a symmetric monoidal category $\cC$ and $X \in \cC$. A \textit{dual object} of $X$ is an object $X^*$ with maps $\ev_X: X^* \otimes X \to \mathbb{1}$ and $\coev_X: \mathbb{1} \to X \otimes X^*$ such that
    \[
    (\mathrm{id}_X\otimes\ev_X)a_{X,X^*,X}
    (\coev_X\otimes\mathrm{id}_X)=\mathrm{id}_X,
    \]
    \[
    (\ev_X\otimes\mathrm{id}_{X^*})a^{-1}_{X^*,X,X^*}
    (\mathrm{id}_{X^*}\otimes\coev_X)=\mathrm{id}_{X^*},
    \]
    where the canonical unit constraints identify the sources and targets with $X$ and $X^*$.
\end{definition}
\begin{definition}
    $\cC$ is \textit{rigid} if all objects admit a dual.
\end{definition}
\begin{remark}
\hfill
\begin{itemize}
    \item If $X^*$ exists, it is unique up to a unique isomorphism compatible with evaluation and coevaluation. 
    \item Rigidity is a property, not a structure.
    \item If $\cC$ is rigid, we get a contravariant functor $X \mapsto X^*$.
    \item In a symmetric category, we get a natural isomorphism $X \cong (X^*)^*$.
    \item Rigidity can be defined for any monoidal category. Without the symmetric braiding, we get ``left'' and ``right'' duals $X^*$ and ${^*X}$, respectively, such that $X \cong {^*(X^*)}$.
\end{itemize}
\end{remark}
\begin{example}
    In $\bVec_k$, we can let $V^* \coloneqq \Hom_k(V, k)$ with
    \begin{align*}
        \ev_V: \hspace{1mm} &V^* \otimes V \to k \\
        &f \otimes v \mapsto f(v)
    \end{align*}
    and
    \begin{align*}
        \coev_V: k & \to V \otimes V^* \\
        1&\mapsto\sum_i v_i\otimes f_i,
    \end{align*}
    where $\{v_i\}$ is a finite basis and $\{f_i\}$ its dual basis. A coevaluation satisfying the triangle identities exists precisely when $V$ is finite dimensional.
    The tensor $\sum_i v_i\otimes f_i$ corresponds to $\mathrm{id}_V$ under $V\otimes V^*\simeq\operatorname{End}_k(V)$, so it is independent of the basis. Conversely, if a coevaluation satisfying the triangle identities has $\coev(1)=\sum_{i=1}^m v_i\otimes f_i$, the first identity gives $v=\sum_i f_i(v)v_i$ for every $v\in V$. Thus the $v_i$ span and $V$ is finite dimensional.
\end{example}
\begin{definition}
    An abelian $k$-linear category $\cC$ is \textit{Artinian} if
    \begin{itemize}
        \item[1)] Hom sets are finite dimensional $k$-vector spaces
        \item[2)] All objects have finite length. That is, for any $X \in \cC$ there is a family of strict inclusions
        \[0=X_0 \subset X_1 \subset \dots \subset X_n = X\]
        such that each $X_i/X_{i-1}$ is simple.
    \end{itemize}
\end{definition}
\begin{definition}
    A \textit{symmetric tensor category} is a category $\cC$ which has
    \begin{itemize}
        \item[1)] A rigid symmetric monoidal structure. 
        \item[2)] An Artinian abelian $k$-linear structure.
    \end{itemize}
    such that
    \begin{itemize}
        \item[1)] $\otimes$ is $k$-linear on Homs (that is, bilinear).
        \item[2)] $\End_{\cC}(\mathbb{1}) \cong k$. Under these hypotheses the unit is simple; see \cite[Section 4.3]{EGNO}.
    \end{itemize}
\end{definition}
\begin{remark}
    Rigidity implies that tensoring with an object has both a left and a right adjoint, given by tensoring with suitable duals. A functor between abelian categories with both adjoints is exact. Hence $\otimes$ is exact in each variable; bilinearity alone would not imply this. See \cite[Proposition 4.2.1]{EGNO}.
\end{remark}
\begin{exercise}
Show that $G$-graded vector spaces are equivalent to right comodules over the group algebra $kG$, by sending $v\in V_g$ to $v\otimes g$.
\end{exercise}
